% Generated draft: requires mathematical and source review.
\documentclass[11pt]{article}
\usepackage[margin=25mm]{geometry}
\usepackage{fontspec}
\setmainfont{Latin Modern Roman}
\usepackage{amsmath,amssymb,mathtools,unicode-math}
\setmathfont{Latin Modern Math}
\usepackage{enumitem,needspace,xurl,hyperref}
\setlength{\parindent}{0pt}
\setlength{\parskip}{5pt}
\setlength{\emergencystretch}{3em}
\newcommand{\N}{\mathbb{N}}
\newcommand{\Z}{\mathbb{Z}}
\newcommand{\Q}{\mathbb{Q}}
\newcommand{\R}{\mathbb{R}}
\newcommand{\C}{\mathbb{C}}
\newcommand{\F}{\mathbb{F}}
\newcommand{\A}{\mathbb{A}}
\newcommand{\PP}{\mathbb{P}}
\newcommand{\E}{\mathbb{E}}
\newcommand{\eps}{\varepsilon}
\newcommand{\dd}{\,\mathrm{d}}
\begin{document}
\section*{UnsolvedMath ID: 30000321}
\textbf{Source problem code:} catalogue entry
\subsection{Definitions and mathematical statement}
Let \(p\) be a prime and \(G=(\Z/p\Z)^r\), with \(r\) the intended nonnegative integer rank. The cleaned source describes \(C\) as a projective \(\Z G\)-chain complex with \(G\)-action. In the usual algebraic meaning this is a chain complex of projective \(\Z G\)-modules and equivariant differentials. The record does not give a geometric realization or a separate definition of fixed points for such a complex.

Retain the source's notation \(C/G\) and
\[
 j^*:H^n(C/G)\longrightarrow H^n(C).
\]
The hypothesis is the equality of subgroups
\[
 \operatorname{im}(j^*)=p^kH^n(C)
 =\{p^k\alpha:\alpha\in H^n(C)\}.
\]
It is not merely an inclusion, and it does not mean that the image has order or index \(p^k\) unless additional information on the cohomology group is given. The coefficient convention, definition of \(C/G\), map \(j\), degree \(n\) and allowable integer \(k\) are missing. For an algebraic complex, coinvariants are one possible meaning of \(C/G\); a geometric orbit space is another and cannot be silently identified with it.

The proposed conclusion is that there exists a subgroup \(H\leq G\) with
\[
 H\cong(\Z/p\Z)^{r-k}
\]
that has a fixed point on \(C\), in the intended geometric sense. This requires \(0\leq k\leq r\) for the displayed rank to be meaningful, but the supplied text does not state that range. If an associated \(G\)-space \(Z\) is intended, the conclusion would mean \(Z^H\ne\varnothing\), equivalently \(\exists x\in Z\ \forall h\in H\;hx=x\). That is a single common fixed point for the subgroup, not separate fixed points for individual elements. No such \(Z\) or realization hypothesis is provided.

An algebraic fixed submodule or fixed subcomplex always contains zero, so its nonemptiness is not a substitute for the requested geometric fixed point. This is an essential formulation gap, not repaired by treating zero as the asserted point. Possible finiteness, boundedness, top-degree or orientation/cohomology hypotheses must also be recovered rather than inferred from the notation \(H^n\).

The literature triage checked 2026-08-21 reports a fixed-point conjecture extending previously proved special geometric cases, with no general proof identified. Those cases and their hypotheses are absent. The open label and reported source normalization remain dated metadata and do not resolve the chain-complex versus geometric-space issue.
\subsection{Short English statement}
The conjecture seeks a common fixed point for a subgroup of rank \(r-k\) from an exact divisibility description of a cohomological pullback image. The source calls the object an algebraic chain complex but does not define its geometric fixed points, quotient or pullback. Those missing structures are indispensable; the zero vector in a fixed subcomplex cannot serve as the intended conclusion.
\subsection{Sources}
\begin{itemize}
\item Local UnsolvedMath registry, permanent ID 30000321, source code catalogue entry, ``Fixed Points from Equivariant Cohomology Divisibility,'' including original assertion, added chain-complex wording and verification metadata.
\item \emph{Cohomology of Finite Groups: Interactions and Applications (2006)}, Oberwolfach Reports, pp.\ 2375--2434, DOI 10.4171/owr/2005/42: \url{https://doi.org/10.4171/owr/2005/42}. This is the sole supplied literature reference, labelled ``original Oberwolfach report.'' The citation uses 2006, while the DOI and proposed-year metadata use 2005.
\item Literature triage checked 2026-08-21; local record updated 2026-09-16. No external source has been visited or independently reviewed here.
\end{itemize}
\end{document}
